\subsubsection{Eigenfunction Expansions}\label{subsec:symmetric-elliptic-eigenfunction-expansions}

When the coefficients of $L(\bm{x},D)$ satisfy the symmetry condition
\begin{equation}\label{eq:symmetric-elliptic-coefficient-condition}
a_{\sigma\gamma}=a_{\gamma\sigma},
\end{equation}
then it is easy to show that $L$ is symmetric. Moreover, by direct calculation we see that for every $u,v\in H_0^k(\Omega)$ we have
\begin{equation}\label{eq:symmetric-elliptic-bilinear-form-condition}
B[u,v]=B[v,u].
\end{equation}


For any $f,g\in L^2(\Omega)$ this gives us\[\begin{aligned}
(G(f),g)_{L^2(\Omega)}&=(\overline G(f),g)_{L^2(\Omega)}\\
&=\dfrac{1}{\widetilde\lambda}B_{\widetilde\lambda}[\overline G(f),\overline G(g)]\\
&=\dfrac{1}{\widetilde\lambda}B_{\widetilde\lambda}[\overline G(g),\overline G(f)]\\
&=(\overline G(g),f)_{L^2(\Omega)}\\
&=(f,G(g))_{L^2(\Omega)}.
\end{aligned}\]So $G$ is self-adjoint. Thus, we can use the Hilbert-Schmidt theorem (Theorem~\ref{thm:hilbert-schmidt-spectral}) to get the following.

\begin{Theorem}
\label{thm:symmetric-elliptic-eigenfunction-expansions}
If $L$ is symmetric, then there is a sequence of real eigenvalues
\begin{equation}\label{eq:symmetric-elliptic-eigenvalue-sequence}
\widetilde\lambda\leq\lambda_1\leq\lambda_2\leq\cdots\leq\lambda_n\leq\cdots
\end{equation}
with no finite accumulation point and $\lim_{i\to\infty}\lambda_i=\infty$, and an orthonormal set of eigenfunctions $\{\phi_i\}_{i=1}^\infty$ such that
\begin{equation}\label{eq:symmetric-elliptic-eigenfunction-equations}
L\phi_i=\lambda_i\phi_i
\end{equation}
(in the weak sense). Furthermore, if $\mu\ne\lambda_i$, $i=1,2,\ldots,\infty$, then for any $f\in L^2(\Omega)$ the unique weak solution of
\begin{equation}\label{eq:symmetric-elliptic-nonhomogeneous-equation}
L(\bm{x},D)u-\mu u=f
\end{equation}
is given by
\begin{equation}\label{eq:symmetric-elliptic-noneigenvalue-solution-expansion}
u=\sum_{i=1}^\infty\dfrac{(\phi_i,f)}{\lambda_i-\mu}\phi_i.
\end{equation}
If $\mu$ is an eigenvalue; i.e., $\mu=\lambda_j$ for $j$ in some index set $J\subset\mathbb{N}$, then \eqref{eq:symmetric-elliptic-nonhomogeneous-equation} is solvable if and only if
\begin{equation}\label{eq:symmetric-elliptic-eigenvalue-compatibility}
(\phi_j,f)=0,\quad j\in J.
\end{equation}
If so, there is a family of solutions given by
\begin{equation}\label{eq:symmetric-elliptic-eigenvalue-solution-family}
u=\sum_{j\in J}\alpha_j\phi_j+\sum_{\mathbb{N}\backslash J}\dfrac{(\phi_i,f)}{\lambda_i-\mu}\phi_i.
\end{equation}
(Here the series \eqref{eq:symmetric-elliptic-noneigenvalue-solution-expansion} and \eqref{eq:symmetric-elliptic-eigenvalue-solution-family} converge in $L^2(\Omega)$.)
\end{Theorem}

The proof is left to the reader.
